\documentclass[11pt]{article}
\usepackage[a4paper,margin=27mm]{geometry}
\usepackage{amsmath,amssymb,amsthm,booktabs,array}
\usepackage[T1]{fontenc}
\usepackage{lmodern}
\usepackage[colorlinks=true,allcolors=blue,breaklinks=true]{hyperref}
\newtheorem{theorem}{Theorem}
\newtheorem{lemma}[theorem]{Lemma}
\newtheorem{corollary}[theorem]{Corollary}
\newcommand{\ds}{\operatorname{d}}
\title{Boundary bounds for Mertens's grid--cylinder conjecture:\\an alternative proof and extensions}
\author{Haobo Ma\\\small ChronoAI Pte Ltd, Singapore
\and Wenlin Zhang\\\small National University of Singapore\\\small ChronoAI Pte Ltd, Singapore}
\date{12 November 2026}
\begin{document}
\maketitle
\begin{abstract}
We give a finite-boundary alternative proof of Mertens's grid--cylinder
growth equality (2024, Section 5, page 19, after equation (22)), which
already follows from Gangloff and Talon's grid results (2019/2020).
The per-vertex counts have the same unrestricted joint limit as both side
lengths tend to infinity. An injective boundary repair gives the finite
comparison $g_{m,n}\le c_{m,n}\le4^n g_{m,n}$.
A block argument proves existence of the grid limit and gives an explicit
finite-rectangle error bound. The same boundary comparison also yields
the common torus limit. Further corollaries give strip error bounds and
extend the comparison to domination polynomials at every positive activity,
including the king graph.
\end{abstract}

\section{The question and its scope}
Mertens \cite[Section 5, p.~19, after (22)]{mertens} conjectures that
the grid and cylinder growth constants agree: ``these two growth rates
are in fact equal.'' We give a finite comparison uniformly over rectangles
of arbitrarily varying aspect ratio. The equality is already a consequence
of Gangloff and Talon \cite[Lemma 4 and Theorem 1]{gangloff}, as explained
below; our focus is an explicit same-size estimate and its extensions.

A set $S$ of vertices of a finite simple graph dominates the graph if
each vertex is in $S$ or has a neighbor in $S$. Write $\ds(H)$ for the
number of all dominating subsets of $H$, counting labelled subsets
without quotienting by symmetries. Let $P_m$ be the path on
$\{0,\ldots,m-1\}$, and let $C_m$, for $m\ge3$, be the cycle obtained
by adding the edge between its endpoints. On the common vertex set
$\{0,\ldots,m-1\}\times\{0,\ldots,n-1\}$ put
\[
 g_{m,n}=\ds(P_m\mathbin\square P_n),\qquad
 c_{m,n}=\ds(C_m\mathbin\square P_n).
\]
Here $\square$ denotes the Cartesian product: an edge changes exactly one
coordinate along an edge of its factor. All grid dimensions are positive;
the cylinder requires $m\ge3$. Natural logarithms are used throughout.

For clarity, let $N_{m,n}$ count restrictions to an $m\times n$ rectangle
of dominating configurations on the infinite square grid. Periodically
repeating a torus dominating set makes such a configuration. Conversely,
surrounding a restriction by a fully selected frame dominates the enlarged
rectangle and retains the restriction. Thus, for $m,n\ge3$,
\[
 g_{m,n}\le c_{m,n}\le t_{m,n}\le N_{m,n}\le g_{m+2,n+2},
 \qquad t_{m,n}=\ds(C_m\square C_n).
\]
The final inequality is the selected-frame construction of
\cite[Lemma 4, p.~91]{gangloff}; its Theorem 1 gives the joint grid
limit. Since $(m+2)(n+2)/(mn)\to1$, the same common limit follows for
cylinders and tori. This is a consequence of that earlier work, rather
than an explicit grid--cylinder statement there. The grid limit is also
stated by Mertens in (15)--(17), and the torus--cylinder equality is his
conclusion on p.~19. Talon \cite[Lemmas 3.6--3.7, pp.~62--63]{talon}
also explicitly treats dimensions tending to infinity at different speeds.
The next proof gives a same-size bound without the added frame.

\begin{theorem}\label{thm:main}
There is a number $\mu\in[1,2]$ such that
\[
 \lim_{m,n\to\infty}g_{m,n}^{1/(mn)}
 =\lim_{m,n\to\infty}c_{m,n}^{1/(mn)}=\mu.
\]
These are unrestricted joint limits: for every $\varepsilon>0$ there is
$N\ge3$ such that both expressions differ from $\mu$ by less than
$\varepsilon$ whenever $m,n\ge N$. Moreover,
\begin{equation}\label{eq:seam}
 g_{m,n}\le c_{m,n}\le4^n g_{m,n}
 \qquad(m\ge3,\ n\ge1).
\end{equation}
\end{theorem}

\section{Two injections}
Adding edges preserves every dominating set, so $\ds(G)\le\ds(H)$
whenever $G\subseteq H$ have the same vertex set. The reverse comparison
can be bounded by the vertices incident to the extra edges.

\begin{lemma}[Boundary repair]\label{lem:repair}
Let $G,H$ be graphs on the same finite vertex set $V$. Suppose $B\subseteq V$
has the property that every $H$-edge incident to a vertex outside $B$
is a $G$-edge. Then
\[
 \ds(H)\le 2^{|B|}\ds(G).
\]
\end{lemma}
\begin{proof}
If $S$ dominates $H$, then $S\cup B$ dominates $G$. Vertices in $B$
are selected, and a vertex outside $B$ retains any neighbor that
dominated it. The map
\[
 S\longmapsto (S\cup B,\ S\cap B)
\]
takes values in the dominating sets of $G$ times the subsets of $B$.
It is injective: the first coordinate determines $S\setminus B$, and
the second determines $S\cap B$. Counting these pairs proves the bound.
\end{proof}

For the cylinder and grid, take $B$ to be the two endpoint columns,
of size $2n$. Outside these columns the graphs have identical neighbors.
Edge inclusion and Lemma~\ref{lem:repair} give \eqref{eq:seam}, including
$m=3$ and $n=1$.

\begin{lemma}[Block injection]\label{lem:block}
For all positive integers $a,b,M,N$,
\begin{equation}\label{eq:block}
 g_{M,N}\ge g_{a,b}^{\lfloor M/a\rfloor\lfloor N/b\rfloor}.
\end{equation}
\end{lemma}
\begin{proof}
Place $q=\lfloor M/a\rfloor$ by $p=\lfloor N/b\rfloor$ disjoint
$a\times b$ rectangular blocks in the larger grid. Choose a dominating
set in each block independently and select every vertex outside the
blocks. Each block vertex is dominated inside its block, and every
remaining vertex is selected. The choices can be recovered by restricting
to each block, so the construction is injective. Edges between blocks
cannot destroy domination. The construction also applies when $q$ or
$p$ is zero, yielding the full vertex set.
\end{proof}

\section{The joint limit and a finite-size estimate}
Define
\[
 f_{m,n}=\frac{\log g_{m,n}}{mn},\qquad
 h=\sup_{a,b\ge1}f_{a,b}.
\]
The full vertex set dominates, and there are $2^{mn}$ subsets altogether.
Thus $0\le f_{m,n}\le\log2$ and $0\le h\le\log2$.
For fixed $a,b$, Lemma~\ref{lem:block} gives
\[
 f_{M,N}\ge
 \frac{a\lfloor M/a\rfloor}{M}
 \frac{b\lfloor N/b\rfloor}{N}f_{a,b}.
\]
Both fractions tend to one as $M,N$ tend jointly to infinity. Consequently
$\liminf f_{M,N}\ge f_{a,b}$ for every $a,b$, hence
$\liminf f_{M,N}\ge h$. Since $f_{M,N}\le h$, the joint limit is $h$.
Equivalently, one can first choose a block with rate within a prescribed
error of $h$ and then choose both larger dimensions so that its covered
area fraction is sufficiently close to one; this gives a single common
threshold on the two dimensions.

Taking logarithms of \eqref{eq:seam} yields
\begin{equation}\label{eq:rate}
 0\le\frac{\log c_{m,n}}{mn}-f_{m,n}
 \le\frac{2\log2}{m}.
\end{equation}
The right side tends to zero at every joint approach to infinity.
Exponentiation proves Theorem~\ref{thm:main}, with $\mu=e^h$.

\begin{corollary}\label{cor:error}
Every positive rectangle gives the rigorous enclosure
\begin{equation}\label{eq:error}
 0\le h-f_{a,b}\le2\log2\left(\frac1a+\frac1b\right).
\end{equation}
\end{corollary}
\begin{proof}
Partition the $qa\times pb$ grid into $qp$ blocks of size $a\times b$,
and delete all edges between blocks to obtain a disjoint union $G$.
Its count is $g_{a,b}^{qp}$. Let $B$ consist of both columns adjoining
each vertical cut and both rows adjoining each horizontal cut. Repeated
rows or columns are counted only once, so
\[
 |B|\le 2(q-1)pb+2(p-1)qa.
\]
Every deleted edge has both endpoints in $B$. Boundary repair gives
\[
 f_{qa,pb}\le f_{a,b}
 +2\log2\left(\frac{q-1}{qa}+\frac{p-1}{pb}\right).
\]
Let $q,p$ tend to infinity and use the joint grid limit. The other
inequality follows from the definition of $h$.
\end{proof}

\begin{corollary}\label{cor:torus}
For $m,n\ge3$, put $t_{m,n}=\ds(C_m\mathbin\square C_n)$. Then
\[
 g_{m,n}\le c_{m,n}\le t_{m,n}\le4^m c_{m,n}
 \le4^{m+n}g_{m,n},\qquad
 \lim_{m,n\to\infty}t_{m,n}^{1/(mn)}=\mu.
\]
\end{corollary}
\begin{proof}
The first two inequalities follow by adding edges. For the next one,
apply boundary repair to the torus and cylinder with $B$ the two
endpoint rows, of size $2m$. The logarithmic difference between torus
and grid is at most $2\log2(1/m+1/n)$, which tends to zero jointly.
\end{proof}

\begin{corollary}\label{cor:strip}
For each fixed $m\ge1$, the limit
$h_m=\lim_{n\to\infty}\log g_{m,n}/(mn)$ exists and
\[
 0\le h-h_m\le \frac{2\log2}{m}.
\]
For $m\ge3$ the cylinder strip limit $\bar h_m$ exists, the torus
has the same strip limit, and
\[
 h_m\le\bar h_m\le h_m+\frac{2\log2}{m}.
\]
In particular both strip limits converge to $h$ as $m\to\infty$.
\end{corollary}
\begin{proof}
For either the grid or cylinder of fixed width $m$, concatenating two
dominating sets gives an injection, so its logarithmic count $a_n$ is
superadditive. To see convergence directly, fix $b$, write $n=qb+r$
with $0\le r<b$, and fill the last $r$ columns. Then $a_n\ge q a_b$.
Since $0\le a_n\le mn\log2$, this proves
$\lim a_n/n=\sup_b a_b/b$. Letting $n\to\infty$ in
\eqref{eq:error} gives the first estimate; taking the same limit in
\eqref{eq:rate} gives the second. For the torus use
$c_{m,n}\le t_{m,n}\le4^m c_{m,n}$: the logarithmic discrepancy per
vertex is at most $2\log2/n$. These facts reprove the strip conclusions
in Mertens (17) and on p.~19, and give an explicit error bound.
\end{proof}

\section{Positive activities and the king graph}
For $z>0$, let $Z_H(z)=\sum_{S\text{ dominates }H}z^{|S|}$ and set
$\kappa_z=1+z^{-1}$. The preceding repair injection gives the weighted
bound
\begin{equation}\label{eq:weighted-repair}
 Z_H(z)\le\kappa_z^{|B|} Z_G(z).
\end{equation}
Indeed, with $T=S\cup B$ and $U=S\cap B$,
$z^{|S|}=z^{|T|}z^{|U|-|B|}$. Summing over the image of the injection,
and then over all pairs $(T,U)$, multiplies $Z_G(z)$ by
$z^{-|B|}(1+z)^{|B|}=\kappa_z^{|B|}$.

\begin{corollary}\label{cor:weighted}
Fix $z>0$ and either the Cartesian grid family or the king graph family
$P_m\boxtimes P_n$. Within the chosen family, the free rectangles,
their one-direction periodic versions, and their two-direction periodic
versions have a common joint logarithmic rate $h(z)$, depending on the
family. The king adjacency
means that distinct vertices differ by at most one in each coordinate;
periodic versions use the corresponding cyclic coordinate distance.
If $F_{a,b}(z)=\log Z_{a,b}(z)/(ab)$ for free rectangles, then
\begin{equation}\label{eq:weighted-error}
 h(z)=\sup_{a,b\ge1}F_{a,b}(z),\qquad
 0\le h(z)-F_{a,b}(z)
 \le 2\log\kappa_z\left(\frac1a+\frac1b\right).
\end{equation}
The finite periodic/free ratios are at most $\kappa_z^{2n}$ for one
seam and $\kappa_z^{2(m+n)}$ for two seams.
\end{corollary}
\begin{proof}
Put $A=MN$, $k=\lfloor M/a\rfloor\lfloor N/b\rfloor$ and
$\theta=kab/A$. Block choices and the selected leftover vertices give
\[
 Z_{M,N}(z)\ge Z_{a,b}(z)^k z^{A-kab},\qquad
 F_{M,N}(z)\ge\theta F_{a,b}(z)+(1-\theta)\log z.
\]
The bounds $z^A\le Z_{M,N}(z)\le(1+z)^A$ ensure that the supremum
of all $F_{a,b}(z)$ is finite. As $M,N\to\infty$, $\theta\to1$;
the same supremum argument as before proves the joint free limit.
For an exact tiling, every deleted Cartesian or diagonal king edge has
both endpoints in the two rows or columns adjacent to a cut.
Apply \eqref{eq:weighted-repair} to those vertices and let the numbers
of blocks tend to infinity. This gives \eqref{eq:weighted-error}.
All extra periodic edges likewise have both endpoints in the seam
columns or rows, so the same repair gives the stated ratios and common
limits. Periodic dimensions are at least three. The block and repair
arguments use only the stated adjacencies, so they apply to both families.
\end{proof}

At $z=1$ this includes existence and finite-rectangle enclosures for the
king-graph rate mentioned after Mertens (17). It does not certify the ten
decimal digits conjectured after (24), or the power-series extrapolation
assumed in (20). In particular the $O(1/m)$ strip bound establishes a rate
bound, not a limiting first-order coefficient as suggested in (19).

\section{Coverage and mathematical evidence}
The common torus--cylinder constant is already described by Mertens
\cite[p.~19]{mertens} using transfer matrices. Corollary~\ref{cor:torus}
provides a direct finite-size comparison. Table~\ref{tab:coverage} records
the nearby statements in the same source.

\begin{table}[ht]
\centering\small
\begin{tabular}{@{}p{.43\linewidth}p{.51\linewidth}@{}}
\toprule
Source item in Mertens (2024) & Coverage\\
\midrule
Section 5, (15)--(17): grid growth & Reproved; prior proof Gangloff--Talon
(2019/2020), Theorem 1; also Mertens (2024).\\[3pt]
Section 5, (17), and following sentence: strip limits and king growth &
Reproved; prior statement Mertens (2024). Corollaries~\ref{cor:strip}
and \ref{cor:weighted} give quantitative extensions.\\[3pt]
Section 5, p.~19, after (22): grid--cylinder equality & Reproved;
consequence of prior results of Gangloff--Talon (2019/2020).
Theorem~\ref{thm:main} and \eqref{eq:seam} give a same-size bound.\\[3pt]
Section 5, p.~19: torus--cylinder equality & Reproved; prior proof
Mertens (2024). Corollaries~\ref{cor:torus} and \ref{cor:strip}
give finite, joint, and strip bounds.\\[3pt]
Section 5, (18)--(22): ratio convergence, first-order behavior,
series and numerical estimates & Not covered as stated; the weaker
$O(1/m)$ strip error is proved.\\[3pt]
Section 5, after (24): ten digits for king graphs & Not covered;
the king-graph numerical conjecture is separate.\\[3pt]
Section 4 and Appendix B: exact polynomials and minimum-set counts;
Appendix A: signature counts & Not covered in general. The accompanying
finite checks include small domination polynomials.\\[3pt]
Positive-activity and periodic king extensions & Not posed; proved in
Corollary~\ref{cor:weighted}.\\[3pt]
OEIS A218354 and A218663: empirical fixed-width recurrences & Not covered;
strip growth existence does not identify those recurrences.\\
\bottomrule
\end{tabular}
\caption{Statements near the posing paragraph and their scope.}
\label{tab:coverage}
\end{table}

The pinned Lean file
\href{https://github.com/the-omega-institute/trureturing/blob/e277c80d045570daf43bd03b4e313f7ce457b4e2/D5/S3/Combinatorics/Graph/Domination/MertensGridGrowth.lean}{\texttt{MertensGridGrowth.lean}}
contains declarations \texttt{finite\_seam\_bound}, \texttt{tiling\_count},
and \texttt{result}, respectively for \eqref{eq:seam} with $m\ge3$,
\eqref{eq:block}, and the unrestricted joint grid--cylinder conclusion.
The same revision's
\href{https://github.com/the-omega-institute/trureturing/blob/e277c80d045570daf43bd03b4e313f7ce457b4e2/D5/S3/Combinatorics/Graph/DominatingSetAverage.lean}{\texttt{DominatingSetAverage.lean}}
defines \texttt{domSets} as all ordinary dominating subsets, supplying
the counting definition used by those declarations.
The entropy, torus, strip, positive-activity, and king-graph extensions
are written proofs. The accompanying \texttt{verify.py} independently
enumerates dominating subsets on small grids, cylinders, tori, and king
graphs, checks repair injections, and tests weighted block and cut
inequalities at $z=1/2,1,2$ using exact rational arithmetic. Its range is
$1\le m,n\le6$ and $mn\le16$, with periodic dimensions at least three.
These finite checks concern the combinatorial constructions; the all-size
limits follow from the proofs above.

\begin{thebibliography}{3}
\bibitem{gangloff}
S.~Gangloff and A.~Talon, Asymptotic growth rate of square grids dominating
sets: a symbolic dynamics approach, \emph{Theoretical Computer Science}
\textbf{821} (2020), 87--101.
\href{https://doi.org/10.1016/j.tcs.2020.03.006}{doi:10.1016/j.tcs.2020.03.006}.
First released 25 June 2019 as
\href{https://arxiv.org/abs/1906.10779}{arXiv:1906.10779}.
\bibitem{talon}
A.~Talon, \emph{Intensive use of computing resources for dominations in
grids and other combinatorial problems}, doctoral thesis, Universit\'e
de Lyon, 2019. Public preprint released 26 February 2020:
\href{https://arxiv.org/abs/2002.11615}{arXiv:2002.11615}.
\bibitem{mertens}
S.~Mertens, Domination polynomials of the grid, the cylinder, the torus,
and the king graph, \emph{Journal of Integer Sequences} \textbf{27}
(2024), Article 24.7.6.
\href{https://cs.uwaterloo.ca/journals/JIS/VOL27/Mertens2/mertens12.pdf}{Full text}.
\end{thebibliography}

\noindent\small\textcopyright\ 2026 Haobo Ma and Wenlin Zhang.
Licensed under Creative Commons Attribution 4.0 International (CC BY 4.0).
\end{document}
